\section{System Architecture}

\subsection{Technology Stack}

\subsubsection{Backend Technologies}

\begin{description}[leftmargin=3cm,style=nextline]
    \item[Framework] Laravel 12.0 (PHP 8.2)
    \item[Authentication] Laravel Sanctum 4.0
    \item[PDF Generation] barryvdh/laravel-dompdf 3.0
    \item[Routing Helper] tightenco/ziggy 2.4 (Inertia route helpers)
    \item[Testing] PHPUnit 11.5, Mockery 1.6
    \item[Code Quality] Laravel Pint 1.18 (PHP CS Fixer)
\end{description}

\subsubsection{Frontend Technologies}

\begin{description}[leftmargin=3cm,style=nextline]
    \item[Framework] React 18.3.1
    \item[Bridge] Inertia.js 2.0.3 (Laravel-React SPA bridge)
    \item[UI Library] Material-UI (MUI) 6.3.0
    \item[Data Grid] @mui/x-data-grid 7.24.1
    \item[Date Pickers] @mui/x-date-pickers 7.24.1
    \item[Charts] Recharts 3.6.0
    \item[CSS Framework] Tailwind CSS 3.4.17
    \item[Build Tool] Vite 6.0.7
    \item[Notifications] react-hot-toast 2.5.1
    \item[Utilities] dayjs 1.11.13, lodash.debounce 4.0.8
\end{description}

\subsubsection{Database \& Infrastructure}

\begin{description}[leftmargin=3cm,style=nextline]
    \item[Database] MySQL 5.7+ / MariaDB 10.3+
    \item[Web Server] Nginx or Apache (Laragon WAMP stack)
    \item[Process Manager] Laravel Queue Worker (database driver for dev, Redis/Horizon for production)
    \item[Scheduler] Laravel Task Scheduler (cron-based)
    \item[Storage] Local filesystem (with symlink to \texttt{public/storage})
\end{description}

\subsection{Backend Architecture}

\subsubsection{Directory Structure}

\begin{lstlisting}[language=bash,label=lst:backend-structure,caption={Laravel Backend Structure}]
app/
├── Console/Commands/          # Artisan commands
├── Domains/Products/          # Domain-driven design modules
├── Events/                    # Event classes
│   ├── BillCancelled.php
│   ├── ProductPriceChanged.php
│   ├── PurchaseReceived.php
│   └── StockAdjusted.php
├── Http/
│   ├── Controllers/          # 15 controllers
│   │   ├── BillsController.php
│   │   ├── ProductsController.php
│   │   ├── PurchasesController.php
│   │   ├── ReplenishmentController.php
│   │   └── ...
│   ├── Middleware/           # Custom middleware
│   └── Requests/             # Form request validation
├── Jobs/                     # Background jobs
│   ├── CheckLowStockJob.php
│   ├── ComputeDailyAnalyticsJob.php
│   ├── ComputeReorderSuggestionsJob.php
│   └── DetectAnomaliesJob.php
├── Listeners/                # Event listeners
│   ├── LogBillCancellation.php
│   ├── LogProductPriceChange.php
│   ├── LogPurchaseReceived.php
│   └── LogStockAdjustment.php
├── Models/                   # Eloquent models (17 models)
│   ├── Product.php
│   ├── Bill.php
│   ├── PurchaseOrder.php
│   ├── ReorderSuggestion.php
│   └── ...
├── Observers/                # Model observers
├── Policies/                 # Authorization policies
│   └── BillPolicy.php
├── Providers/
│   └── AppServiceProvider.php
└── Services/                 # Business logic services
    ├── Analytics/
    │   ├── SalesAnalyticsService.php
    │   └── InventoryAnalyticsService.php
    ├── Forecasting/
    │   └── DemandForecastService.php
    └── NotificationService.php
\end{lstlisting}

\subsubsection{Routing Architecture}

The application uses a single web routes file (\texttt{routes/web.php}) with middleware-based route groups:

\begin{lstlisting}[language=PHP,label=lst:routing,caption={Route Structure Example}]
// Public routes
Route::middleware('guest')->group(function () {
    Route::get('login', [AuthenticatedSessionController::class, 'create']);
    Route::post('login', [AuthenticatedSessionController::class, 'store']);
});

// Protected routes
Route::middleware('auth')->group(function () {
    Route::get('/', [DashboardController::class, 'index']);
    
    // Resource routes with additional actions
    Route::resource('products', ProductsController::class);
    Route::post('products/{product}/adjust-stock', 
        [ProductsController::class, 'adjustStock']);
    
    // Nested route groups
    Route::prefix('replenishment')->name('replenishment.')->group(function () {
        Route::get('/', [ReplenishmentController::class, 'index']);
        Route::post('/approve', [ReplenishmentController::class, 'approve']);
        // ...
    });
});
\end{lstlisting}

\subsubsection{Controller Design Pattern}

Controllers follow a consistent pattern:

\begin{enumerate}
    \item \textbf{Validation:} Form Request classes or inline \texttt{\$request->validate()}
    \item \textbf{Authorization:} Policy checks via \texttt{authorize()} or middleware
    \item \textbf{Business Logic:} Delegated to Service classes or Eloquent models
    \item \textbf{Response:} \texttt{Inertia::render()} for pages, JSON for AJAX
\end{enumerate}

\textbf{Example:} \texttt{BillsController::store()}

\begin{lstlisting}[language=PHP,label=lst:controller-pattern]
public function store(Request $request) {
    $validated = $request->validate([
        'worker_id' => 'nullable|exists:workers,id',
        'items' => 'required|array|min:1',
        'items.*.product_id' => 'required|exists:products,id',
        // ...
    ]);
    
    DB::beginTransaction();
    try {
        $bill = Bill::create([/* ... */]);
        foreach ($validated['items'] as $item) {
            $bill->items()->create($item);
            $product->decrement('quantity', $item['quantity']);
        }
        DB::commit();
        return redirect()->route('bills.show', $bill);
    } catch (\Exception $e) {
        DB::rollBack();
        throw $e;
    }
}
\end{lstlisting}

\subsubsection{Service Layer}

Complex business logic resides in dedicated service classes:

\begin{itemize}
    \item \textbf{SalesAnalyticsService:} Aggregates sales data, computes KPIs
    \item \textbf{InventoryAnalyticsService:} Calculates inventory turnover, stock aging
    \item \textbf{DemandForecastService:} Implements forecasting algorithms (moving average, exponential smoothing)
    \item \textbf{NotificationService:} Creates and dispatches notifications
\end{itemize}

\subsubsection{Job Queue Architecture}

Background jobs implement \texttt{ShouldQueue} interface and are processed by Laravel Queue Workers:

\begin{description}[style=nextline]
    \item[CheckLowStockJob] Runs daily at 06:00, creates low-stock notifications
    \item[ComputeDailyAnalyticsJob] Runs daily at 02:00, aggregates previous day's sales/inventory
    \item[ComputeReorderSuggestionsJob] Runs daily at 02:30, generates replenishment recommendations
    \item[DetectAnomaliesJob] Runs every 6 hours, identifies unusual patterns
\end{description}

Job execution is tracked via \texttt{JobRun} model with status (\texttt{running}, \texttt{completed}, \texttt{failed}), duration, and metadata.

\subsubsection{Event-Driven Architecture}

Critical operations fire Laravel Events handled by Listeners:

\begin{center}
\begin{tabular}{@{}ll@{}}
\toprule
\textbf{Event} & \textbf{Listener} \\
\midrule
\texttt{BillCancelled} & \texttt{LogBillCancellation} \\
\texttt{ProductPriceChanged} & \texttt{LogProductPriceChange} \\
\texttt{StockAdjusted} & \texttt{LogStockAdjustment} \\
\texttt{PurchaseReceived} & \texttt{LogPurchaseReceived} \\
\bottomrule
\end{tabular}
\end{center}

Listeners write to \texttt{audit\_logs} table and can trigger notifications or additional business logic.

\subsection{Frontend Architecture}

\subsubsection{Inertia.js Bridge}

Inertia.js eliminates the need for a REST API by allowing controllers to return React components directly:

\begin{lstlisting}[language=PHP,label=lst:inertia-backend]
// Laravel Controller
return Inertia::render('Products/Index', [
    'products' => $products,
    'filters' => $filters,
]);
\end{lstlisting}

\begin{lstlisting}[language=JavaScript,label=lst:inertia-frontend]
// React Component (Products/Index.jsx)
export default function Index({ products, filters }) {
    return (
        <AuthenticatedLayout>
            <DataGrid rows={products.data} />
        </AuthenticatedLayout>
    );
}
\end{lstlisting}

Inertia provides:
\begin{itemize}
    \item Server-side routing with client-side rendering
    \item Automatic CSRF protection
    \item Form helpers (\texttt{useForm} hook)
    \item Ziggy route helpers in JavaScript
\end{itemize}

\subsubsection{Component Structure}

\begin{lstlisting}[language=bash,label=lst:frontend-structure,caption={React Frontend Structure}]
resources/js/
├── Components/              # Shared components
│   ├── Layout/
│   │   └── AuthenticatedLayout.jsx  # Main app shell
│   ├── UI/
│   │   ├── DataTable.jsx
│   │   ├── KPICard.jsx
│   │   ├── FilterPanel.jsx
│   │   └── ...
│   └── Forms/
│       ├── ProductForm.jsx
│       └── BillForm.jsx
├── Pages/                  # Inertia page components
│   ├── Dashboard.jsx
│   ├── Bills/
│   │   ├── Index.jsx
│   │   ├── Create.jsx
│   │   └── Show.jsx
│   ├── Products/
│   │   ├── Index.jsx
│   │   ├── Create.jsx
│   │   └── Edit.jsx
│   └── ...
├── hooks/                  # Custom React hooks
│   ├── useTheme.js
│   ├── useFilters.js
│   └── useDebounce.js
├── theme/                  # MUI theme configuration
│   ├── tokens.js           # Design tokens
│   ├── lightTheme.js
│   └── darkTheme.js
└── app.jsx                 # Inertia app entry point
\end{lstlisting}

\subsubsection{State Management}

The application uses a hybrid state management approach:

\begin{itemize}
    \item \textbf{Server State:} Managed by Inertia.js (automatic synchronization on navigation)
    \item \textbf{Local UI State:} React \texttt{useState} and \texttt{useReducer} hooks
    \item \textbf{Global Settings:} React Context (\texttt{ThemeContext}, \texttt{LocaleContext})
    \item \textbf{Persistent State:} \texttt{localStorage} for theme and language preferences
\end{itemize}

No Redux or external state library is required due to Inertia's server-driven model.

\subsubsection{UI Design System}

Custom MUI theme based on design tokens (\texttt{theme/tokens.js}):

\begin{lstlisting}[language=JavaScript,label=lst:theme-tokens]
export const tokens = {
    colors: {
        primary: { main: '#2563EB', light: '#3B82F6', dark: '#1E40AF' },
        secondary: { main: '#7C3AED', light: '#8B5CF6', dark: '#6D28D9' },
        success: { main: '#10B981' },
        warning: { main: '#F59E0B' },
        error: { main: '#EF4444' },
    },
    typography: {
        fontFamily: "'Inter', 'Roboto', sans-serif",
        h1: { fontSize: '2.5rem', fontWeight: 700 },
        body1: { fontSize: '1rem', lineHeight: 1.5 },
    },
    spacing: {
        xs: 4, sm: 8, md: 16, lg: 24, xl: 32,
    },
};
\end{lstlisting}

Light and dark themes extend these tokens with mode-specific background and text colors.

\subsection{Database Architecture}

\subsubsection{Entity-Relationship Overview}

The database consists of 17 primary tables organized into functional groups:

\begin{center}
\begin{longtable}{@{}p{4cm}p{9cm}@{}}
\caption{Database Tables by Category}\\
\toprule
\textbf{Category} & \textbf{Tables} \\
\midrule
\endfirsthead
\multicolumn{2}{c}{\textit{(continued from previous page)}}\\
\toprule
\textbf{Category} & \textbf{Tables} \\
\midrule
\endhead
\midrule
\multicolumn{2}{r}{\textit{(continued on next page)}}\\
\endfoot
\bottomrule
\endlastfoot
Core Inventory & \texttt{products}, \texttt{categories}, \texttt{suppliers}, \texttt{product\_supplier} (pivot) \\
Sales & \texttt{bills}, \texttt{bill\_items} \\
Procurement & \texttt{purchase\_orders}, \texttt{purchase\_order\_items} \\
Analytics & \texttt{analytics\_snapshots}, \texttt{inventory\_movements}, \texttt{reorder\_suggestions}, \texttt{anomaly\_findings} \\
Administration & \texttt{users}, \texttt{workers}, \texttt{settings}, \texttt{audit\_logs}, \texttt{notifications}, \texttt{job\_runs} \\
Framework & \texttt{cache}, \texttt{jobs}, \texttt{failed\_jobs}, \texttt{sessions} \\
\end{longtable}
\end{center}

\subsubsection{Core Table Relationships}

\begin{figure}[h]
\centering
\begin{tikzpicture}[
    node distance=2cm,
    entity/.style={rectangle, draw, fill=blue!10, text width=2.5cm, text centered, minimum height=1cm},
    relationship/.style={diamond, draw, fill=orange!20, text width=1.5cm, text centered, aspect=2},
    arrow/.style={->, >=stealth, thick}
]

% Entities
\node[entity] (product) {Product};
\node[entity, below left=of product] (category) {Category};
\node[entity, below right=of product] (supplier) {Supplier};
\node[entity, right=3cm of product] (bill) {Bill};
\node[entity, below=of bill] (billitem) {BillItem};
\node[entity, above right=of product] (po) {PurchaseOrder};
\node[entity, right=of po] (poitem) {POItem};

% Relationships
\draw[arrow] (product) -- (category) node[midway, above, sloped] {belongs to};
\draw[arrow] (product) -- (supplier) node[midway, above, sloped] {has many};
\draw[arrow] (billitem) -- (product);
\draw[arrow] (billitem) -- (bill);
\draw[arrow] (poitem) -- (product);
\draw[arrow] (poitem) -- (po);
\draw[arrow] (po) -- (supplier);

\end{tikzpicture}
\caption{Simplified Entity Relationship Diagram}
\label{fig:erd}
\end{figure}

\subsubsection{Key Schema Details}

\textbf{Products Table}

\begin{lstlisting}[language=SQL]
CREATE TABLE products (
    id BIGINT UNSIGNED PRIMARY KEY AUTO_INCREMENT,
    name VARCHAR(255) NOT NULL,
    name_ar VARCHAR(255),
    sku VARCHAR(100) UNIQUE,
    barcode VARCHAR(100) UNIQUE,
    description TEXT,
    category_id BIGINT UNSIGNED,
    purchase_price DECIMAL(10,2) NOT NULL,
    selling_price DECIMAL(10,2) NOT NULL,
    quantity DECIMAL(10,2) DEFAULT 0,
    unit VARCHAR(50),
    min_stock DECIMAL(10,2) DEFAULT 0,
    location VARCHAR(100),
    image VARCHAR(255),
    is_active BOOLEAN DEFAULT TRUE,
    created_at TIMESTAMP,
    updated_at TIMESTAMP,
    INDEX idx_barcode (barcode),
    INDEX idx_sku (sku),
    INDEX idx_category (category_id),
    INDEX idx_active_stock (is_active, quantity)
);
\end{lstlisting}

\textbf{Bills \& Bill Items}

\begin{lstlisting}[language=SQL]
CREATE TABLE bills (
    id BIGINT UNSIGNED PRIMARY KEY AUTO_INCREMENT,
    bill_number VARCHAR(50) UNIQUE,
    user_id BIGINT UNSIGNED NOT NULL,
    worker_id BIGINT UNSIGNED,
    customer_name VARCHAR(255),
    customer_phone VARCHAR(20),
    subtotal DECIMAL(10,2) DEFAULT 0,
    discount DECIMAL(10,2) DEFAULT 0,
    tax DECIMAL(10,2) DEFAULT 0,
    total DECIMAL(10,2) DEFAULT 0,
    payment_method ENUM('cash','card','check','credit','other'),
    status ENUM('completed','cancelled') DEFAULT 'completed',
    notes TEXT,
    created_at TIMESTAMP,
    updated_at TIMESTAMP,
    INDEX idx_bill_number (bill_number),
    INDEX idx_created_at (created_at),
    INDEX idx_status (status)
);

CREATE TABLE bill_items (
    id BIGINT UNSIGNED PRIMARY KEY AUTO_INCREMENT,
    bill_id BIGINT UNSIGNED NOT NULL,
    product_id BIGINT UNSIGNED NOT NULL,
    product_name VARCHAR(255),
    quantity DECIMAL(10,2) NOT NULL,
    unit_price DECIMAL(10,2) NOT NULL,
    discount DECIMAL(10,2) DEFAULT 0,
    total DECIMAL(10,2) NOT NULL,
    FOREIGN KEY (bill_id) REFERENCES bills(id) ON DELETE CASCADE
);
\end{lstlisting}

\textbf{Reorder Suggestions}

\begin{lstlisting}[language=SQL]
CREATE TABLE reorder_suggestions (
    id BIGINT UNSIGNED PRIMARY KEY AUTO_INCREMENT,
    product_id BIGINT UNSIGNED NOT NULL,
    suggested_supplier_id BIGINT UNSIGNED,
    recommended_qty DECIMAL(10,2) NOT NULL,
    current_stock DECIMAL(10,2) NOT NULL,
    reorder_point DECIMAL(10,2) NOT NULL,
    safety_stock DECIMAL(10,2),
    avg_daily_demand DECIMAL(10,2),
    projected_stockout_date DATE,
    urgency ENUM('critical','high','medium','low'),
    confidence DECIMAL(5,2) DEFAULT 0,
    reason_json JSON,
    forecast_data JSON,
    status ENUM('pending','approved','dismissed','ordered'),
    created_at TIMESTAMP,
    updated_at TIMESTAMP,
    INDEX idx_status_urgency (status, urgency),
    INDEX idx_product (product_id)
);
\end{lstlisting}

\subsubsection{Indexing Strategy}

Performance-critical indexes include:

\begin{itemize}
    \item \textbf{products:} \texttt{barcode}, \texttt{sku}, \texttt{category\_id}, composite \texttt{(is\_active, quantity)}
    \item \textbf{bills:} \texttt{bill\_number}, \texttt{created\_at}, \texttt{status}
    \item \textbf{inventory\_movements:} composite \texttt{(product\_id, type, created\_at)}
    \item \textbf{reorder\_suggestions:} composite \texttt{(status, urgency)}
    \item \textbf{audit\_logs:} composite \texttt{(event\_type, created\_at)}
\end{itemize}

A dedicated migration (\texttt{2024\_01\_01\_100000\_add\_dashboard\_performance\_indexes.php}) adds additional indexes to optimize dashboard queries.
